\worldchapter{Der Kampf}{combat}
\index{Kampf}
\label{ch:combat}

\begin{multicols}{2}

\section{Beginn des Kampfes}
\index{Beginn des Kampfes}

Sobald ein Konflikt entsteht, wird die Zeit eingefroren und die Gruppe bestimmt die Reihenfolge, in der die Teilnehmer des Kampfes agieren.

\subsection{Initiative}
\index{Initiative}

Dazu würfelt jeder Teilnehmer auf W6 entsprechend dem \textit{Schnelligkeitswert}. Auch hier gilt die Regel der explodierenden Würfel. Die Würfelergebnisse werden addiert. Der Teilnehmer mit dem höchsten Ergebnis beginnt den Kampf, alle anderen folgen in der Reihenfolge ihres Ergebnisses.

\begin{pexample}[Beispiel]
Eine Schurkin mit Schnelligkeit 4 wirft auf ihre Schnelligkeit und erhält 4, 5, 1 und 17. Ihr Ergebnis ist 27.
\end{pexample}

Haben zwei Teilnehmer das gleiche Ergebnis, entscheidet zuerst der Wert \textit{Schnelligkeit} und wenn auch dieser gleich ist, der Wert \textit{Geschick}.

\subsection{Blitzreaktion}
\index{Blitzreaktion}

Vor Beginn des Kampfes macht jeder Teilnehmer eine Probe auf seine \textit{Auffassungsgabe}. Diese Probe symbolisiert die Fähigkeit des Charakters, blitzschnell auf neue Bedrohungen zu reagieren. Bei Erfolg erhält der Teilnehmer eine ``Blitzreaktion'', die es ihm ermöglicht, noch vor Beginn seiner ersten Kampfrunde zu reagieren. Diese Blitzreaktion gilt als normale Aktion (siehe unten), darf aber nur für Reaktionen verwendet werden.

Misslingt diese Probe, so erhält der Teilnehmer erst zu Beginn seiner ersten Kampfrunde seine Aktionen, und kann vorher nicht handeln.

Sobald der Spieler regulär an der Reihe ist, werden seine Aktionen aufgefrischt und die Blitzreaktion verfällt.

\section{Ablauf der Runden}
\index{Kampfrunde!Ablauf}

Der Kampf ist in \textit{Kampfrunden} aufgeteilt. Diese haben folgenden Ablauf:
\begin{itemize}
    \item Rundenbeginn
    \item Für jeden Teilnehmer wird in Reihenfolge der Initiative die ``Spielerkampfrunde'' durchgeführt
    \begin{itemize}
        \item Beginn der Spielerkampfrunde
        \item Die Aktionen des Spielers werden aufgefrischt
        \item Der Spieler führt seine Aktionen aus
        \item Ende der Spielerkampfrunde
    \end{itemize}
    \item Rundenende
\end{itemize}

Der Teilnehmer kann für jede seiner verfügbaren Aktionen eine Handlung ausführen (s. \textit{Akteur und Aktionen}) oder die Aktion für eine Reaktion in der Gegnerrunde aufheben (s. \textit{Reaktionen}).

Hat der letzte Teilnehmer agiert, beginnt die nächste \textit{Kampfrunde} mit dem ersten Teilnehmer.

Sowohl ``Rundenbeginn'' als auch ``Rundenende'' sind Phasen, in denen Reaktionen ausgeführt werden können. Dazu müssen die Teilnehmer Aktionen aufheben und den letzten \textit{Akteur} direkt wahrnehmen. Einige Effekte, wie z.B. Zaubersprüche, können auch in diesen Phasen ausgeführt werden.

\subsection{Akteur und Aktionen}
\index{Akteur und Aktionen}

Wenn ein Teilnehmer an der Reihe ist, ist er der \textit{Akteur}. Der \textit{Akteur} ist der Teilnehmer, der aktiv handelt und seine Aktionen nach Belieben einsetzen oder zurückhalten kann.

Zu Beginn der Spielerkampfrunde werden die \textit{Aktionen} des Spielers aufgefrischt. Die Anzahl der Aktionen des Charakters ergibt sich aus den Charakterschablonen. Der Basiswert für alle Charaktere ist 2.

``Auffrischen'' bedeutet also, dass alle Aktionen wieder zur Verfügung stehen. Hat der Teilnehmer bereits Aktionen verbraucht, z.B. durch Handlungen in der vorherigen Kampfrunde, setzt er seine verfügbaren Aktionen auf das Maximum zurück.

Sind die Aktionen aufgefrischt, kann der Teilnehmer im Kampf handeln. Dazu führt er nacheinander Handlungen aus, wobei jede Handlung eine gewisse Zahl an \textit{Aktionen} in Anspruch nimmt. Handlungen können u.~A. folgende sein:
\begin{itemize}
    \item Mit einer Waffe \textbf{angreifen}
    \item Mit einer Waffe oder einem Gegenstand \textbf{parieren}
    \item Eine Waffe \textbf{nachladen}
    \item Einen Gegenstand \textbf{benutzen}
    \item Einem Nahkampfangriff \textbf{ausweichen}
    \item Mit einer Nahkampfwaffe oder beim Einzelschuss \textbf{zielen}
    \item Eine beliebige Aktion ausführen (s.U.)
    \item \textbf{Hinhocken} oder auf den Boden legen (Der ``In der Hocke'' Status Effekt ist aktiv, siehe \cref{ch:wounds})
    \item \textbf{Aufstehen}
    \item \textit{Schnelligkeit} + 1 Meter \textbf{gehen} (und dabei eine weitere Handlung ausführen, ohne dass sie eine Aktion verbraucht, jedoch ist der Mindestwurf um 1 erhöht)
    \item \textit{Schnelligkeit} + 5 Meter \textbf{rennen}
    \item \textit{Schnelligkeit} / 2 + 1 Meter \textbf{robben} (aufgerundet), der Charakter muss \textit{In der Hocke} sein. (siehe \cref{ch:wounds}: Zustände des Charakters)
\end{itemize}

Aktionen sollten nicht zusammen, sondern immer nacheinander ausgeführt werden, da es mögliche Reaktionen gibt.

\begin{center}
    \includegraphics[width=\columnwidth]{images/chapters/combat/fighter-\wbidentifier.png}
    \label{fig:fighter}
\end{center}

\subsubsection{Beliebige Aktionen}
\index{Beliebige Aktionen}

Ein Charakter kann auch beliebige Aktionen ausführen, die nicht auf der Liste stehen. In diesem Fall muss die Spielleiterin entscheiden, ob die Handlung eine oder mehrere Aktionen erfordert. Eine Aktion, die nicht auf der Liste steht, sollte normalerweise eine Aktion erfordern. Das kann alles Mögliche sein, z.B. das Anzünden einer Pfeife, das Zusammenschlagen der Köpfe zweier Gegner oder das Werfen eines Gegenstandes. Die Spielleiterin entscheidet, welcher Wurf erforderlich ist.

\subsection{Reaktionen}
\index{Reaktionen}

Wenn ein \textit{Akteur} im Kampf agiert, können alle Teilnehmer, die den \textit{Akteur} direkt wahrnehmen, auf diese Aktion reagieren.

Um auf eine Aktion reagieren zu können, müssen folgende Bedingungen erfüllt sein:
\begin{itemize}
    \item Der reagierende Teilnehmer muss den \textit{Akteur} direkt wahrnehmen, d.h. er muss dessen Aktion hören, sehen oder auf andere Weise zur Kenntnis nehmen.
    \item Der reagierende Teilnehmer hat noch unverbrauchte \textit{Aktionen}.
\end{itemize}

Die Reaktion wird unmittelbar nach der Aktion des \textit{Akteurs} angesagt und ausgeführt, findet aber in der Spielzeit davor statt. Auf eine \textit{Aktion} kann immer nur eine \textit{Reaktion} eines Teilnehmers folgen. Es können aber beliebig viele Teilnehmer auf den \textit{Akteur} reagieren, wenn sie seine \textit{Aktion} wahrnehmen. In der Praxis bedeutet dies, dass der reagierende Teilnehmer seine Reaktion ansagt, nachdem der \textit{Akteur} seine Aktion ausgeführt und eventuell gewürfelt hat. Dies kann von Situation zu Situation variieren.

Wenn mehrere Teilnehmer auf eine Aktion reagieren, wird die Reihenfolge der Reaktionen durch die Initiative bestimmt. Der Teilnehmer mit der höchsten Initiative reagiert zuerst, gefolgt von den anderen Teilnehmern in absteigender Reihenfolge der Initiative.

Jede \textit{Reaktion} reduziert die verfügbaren \textit{Aktionen} des reagierenden Teilnehmers um eins.

\begin{pexample}[Beispiel]
Hagen ist in einen Kampf mit einem Räuber verwickelt. Hagen hat in seiner Kampfrunde zugeschlagen, sich aber eine Aktion aufgehoben, um reagieren zu können.
Die Kampfrunde des Räubers beginnt. Der Räuber schlägt zu. Die Spielleiterin würfelt mit vier Würfeln drei Erfolge, also drei Treffer.
Hagens Spieler entscheidet, dass Hagen mit einer \textit{Schildparade} reagieren soll. Er sagt die Reaktion auf den Angriff des Räubers an, nachdem die Spielleiterin den Angriff ausgeführt hat. Er kann dies tun, da er noch eine Aktion übrig hat und den Angriff des Gegners direkt wahrnimmt.
Die Reaktion findet nun im Spiel vor dem Angriff des Räubers statt. Durch die Regel der Schildparade erhält Hagen einen Deckungswurf von 5+ für seinen Rundschild. Er würfelt für jeden der drei Treffer des Räubers. Zweimal würfelt er eine 5 und verhindert damit zwei Treffer. Der dritte Treffer trifft ihn.
\end{pexample}

\section{Aktionen durch Bonuswürfel}
\index{Aktionen!durch Bonuswürfel}\index{Bonuswürfel!Aktionen}

Bonus- und Schicksalswürfel können im Kampf verwendet werden, um Aktionen zu erlangen oder zu stehlen.

Um eine zusätzliche Aktion zu erlangen, kann ein \textit{Bonuswürfel} abgezogen werden. Die zusätzliche Aktion steht sofort zur Verfügung, auch für eine Reaktion.

Wird ein \textit{Schicksalswürfel} ausgegeben, kann einem Gegner eine Aktion gestohlen werden. Diese Aktion steht dem Gegner in dieser Runde nicht mehr zur Verfügung. Hat der Gegner eine Aktion für eine Reaktion aufgehoben, wird diese Aktion entfernt. Der Spieler, der den Schicksalswürfel ausgegeben hat, hat die Aktion sofort zur Verfügung, auch als Reaktion.

Das Ausgeben von Würfeln für Aktionen nimmt selbst keine Aktion in Anspruch.

\section{Ablauf eines Angriffs}
\index{Angriff!Ablauf}

Angriffe mit Waffen werden im Nahkampf und im Fernkampf genau gleich gehandhabt. Der einzige Unterschied ist, dass Angriffe mit Nahkampfwaffen auf die Fertigkeit \textit{Nahkampf} geworfen werden, Angriffe mit Schusswaffen auf die Fertigkeit \textit{Schießen} und Angriffe mit Wurfwaffen auf die Fertigkeit \textit{Werfen}.

Ein Angriff hat folgende Phasen:
\begin{itemize}
    \item \textbf{Der Trefferwurf} ermittelt, wie viele Treffer ein Charakter im Angriff mit einer Waffe erreicht. Hier wird auf die jeweilige Waffenfertigkeit gewürfelt, und zwischen \textit{Kritischen Treffern} und \textit{Treffern} unterschieden.
    \item Der \textbf{Deckungswurf} steht dem Angegriffenen zur Verfügung, wenn er Deckung hat. Hier ist es bereits vor dem Einschlag der Treffer in die Rüstung möglich, Schaden abzuwenden. Schilde können Deckung geben.
    \item Das Umwandeln von \textit{Treffern} in \textit{Wunden} unter der Berücksichtigung von \textit{Schutz}, \textit{Durchschlag} und \textit{kritischen Treffern}.
\end{itemize}

\subsection{Der Trefferwurf}
\index{Trefferwurf}

Um einen Angriff durchzuführen wird ein Wurf mit der entsprechenden Fertigkeit gemacht (Schießen, Nahkampf oder Werfen). Die Anzahl der Würfel wird um das Schadenspotential der Waffe erhöht. Der \textit{Mindestwurf} dieses Wurfs entspricht dem \textit{Mindestwurf} des Charakters.

\wbif{tirakan}{}{
Hierbei ist ein möglicher \textit{Rückstossmalus} zu berücksichtigen, wenn der Charakter bereits in derselben Kampfrunde geschossen hat.
}

Die Anzahl der Würfel entspricht zunächst dem jeweiligen Fertigkeitswert des Charakters (Schießen, Nahkampf, Werfen) zuzüglich des \textit{Schadenspotentials} der Waffe.

\wbif{tirakan}{}{
Der Trefferwurf kann auch durch weitere Umstände verändert werden. Unterschiedliche Angriffsmodi und Schüsse bei falscher Distanz sorgen u.U. für eine Änderung der verfügbaren Würfel.
}

Jeder Erfolg verursacht einen \textit{Treffer} beim Ziel des Angriffs. Wie das Ziel Schaden verhindern kann, ist unter \textit{Wunden und Durchschlag} und \textit{Deckung} beschrieben.

\wbif{tirakan}{}{
\subsubsection{Rückstoss}
\index{Rückstoss}

Automatische Waffen verursachen beim Angriff in der Regel einen \textit{Rückstoss}, welcher das erneute Anvisieren eines Ziels bei einem direkt folgenden Angriff erschweren.

Folgt \textit{innerhalb einer Kampfrunde} auf einen Angriff mit einer Schusswaffe \textit{direkt} ein weiterer Angriff vom selben Charakter, so sind der Mindestwurf und die Schwelle für kritische Treffer um 2 erhöht. Dieser Malus erhöht sich für jeden darauffolgenden Angriff in derselben Kampfrunde. Ein dritter Angriff hat also einen Malus von +4 auf den Mindestwurf und die Schwelle für kritische Treffer.

Der Rückstoss kann verhindert werden, wenn in einer Kampfrunde z.B. eine andere Aktion zwischen zwei Angriffe eingefügt wird. So kommt der Rückstoss zum Beispiel mit Bögen nicht zum Tragen, da zwischen den Angriffen ein neuer Pfeil auf die Sehne gelegt werden muss.

Waffen können eine Rückstosskompensation haben. Dieser Wert senkt den Malus pro Angriff. So wird der Mindestwurf bei einem folgenden Angriff mit einer Waffe mit Rückstosskompensation 1 nur um 1 angehoben. Eine Rückstosskompensation von 2 sorgt dafür, dass Rückstoss für die Waffe nicht mehr relevant ist.

Rückstoss wird nicht über Kampfrunden hinweg berücksichtigt, nur innerhalb einer Kampfrunde.
}

\subsection{Kritische Treffer}
\index{Kritische Treffer}

Beim Trefferwurf verursachte \textit{Treffer} werden zu \textit{kritischen Treffern}, wenn sie beim Wurf den Wert 11 erreichen. Das entspricht einem ``weiter geworfenen'' \textit{Explodierenden Würfel}, der danach erneut ein Resultat von 5+ zeigt. Veränderungen des \textit{Mindestwurfs} des Charakters werden hier nicht angewandt.

\wbif{tirakan}{}{
Kritische Treffer können nur von Nahkampf- und Einzelschussangriffen sowie Wurfwaffen verursacht werden, niemals durch Feuerstoß Angriffe.
}

Werden beim Angriff kritische Treffer erreicht, so werden diese getrennt von den normalen Treffern angesagt. Ein Einzelschuss eines Bogens könnte so zu dem Ergebnis ``2 Crits, 3 normale Treffer'' führen.

Kritische Treffer werden wie normale Treffer behandelt, durchdringen jedoch immer normale Rüstung. Für kritische Treffer kann nur Schutz vom Typ ``Schutz vor kritischen Treffern'' schützen, alle anderen Schutzarten können keine kritischen Treffer verhindern.

Kommt es zu einem \textit{Deckungswurf}, so müssen kritische Treffer von normalen Treffern getrennt behandelt werden. Der Angegriffene wirft also zweimal auf seine Deckung, einmal für die Anzahl der kritischen Treffer, und einmal für die Anzahl der normalen Treffer.

\begin{pexample}[Beispiel]
Der Söldner Maragas würfelt und erhält 4, 5, 5 und 14. Er hat also 2 normale Treffer und einen kritischen Treffer. Der kritische Treffer durchschlägt die Rüstung, die normalen Treffer können durch den Schutz des Angegriffenen reduziert werden.
\end{example}

\subsubsection{Megakritische Treffer}
\index{Megakritische Treffer}

Kommt es zu \textit{kritischen Treffern}, so kann bei den \textit{Explodierenden Würfeln} weiter als 11 geworfen werden. Es wird so lange weiter geworfen, bis keine 6 auf dem jeweiligen Würfel mehr erreicht wird.

Erreicht ein Würfel nach dem zweiten Wurf erneut eine 5, so handelt es sich um einen \textit{Megakritischen Treffer}. Diese Treffer werden wie kritische Treffer behandelt, verursachen jedoch eine zusätzliche Wunde, wenn sie nicht verhindert werden.

Für jedes Weiterwürfeln mit einer 5+ wird hierbei die Zahl der Wunden erhöht. Ein megakritischer Treffer kann also sehr viele Wunden verursachen. Aus der Regel der 5+ ergeben sich folgende Grenzen für Wunden:
\begin{itemize}
    \item \textbf{Wurf 5+}: normaler Treffer
    \item \textbf{Wurf 11+}: kritischer Treffer - ignoriert Rüstung
    \item \textbf{Wurf 17+}: megakritischer Treffer - ignoriert Rüstung, +1 Wunde
    \item \textbf{Wurf 23+}: megakritischer Treffer - ignoriert Rüstung, +2 Wunden
    \item \textbf{Wurf 29+}: megakritischer Treffer - ignoriert Rüstung, +3 Wunden
\end{itemize}
Und so weiter.

\subsubsection{Zielen}
\index{Zielen}

\wbif{tirakan}{
Mit einer Waffe kann sowohl im Nah- als auch im Fernkampf gezielt werden.
}{
Bei Nahkampfwaffen und beim Einzelschuss ist es möglich mit der Waffe zu zielen. Im Feuerstoß Modus ist dies nicht möglich.
}

Der Charakter kann Aktionen investieren, um sein Ziel genauer anzuvisieren. Für jeweils 1 Aktion wird hierbei für den nächsten Angriff die Grenze für kritische Treffer um 2 verringert. Dieser Bonus auf kritische Treffer darf den Wahrnehmungswert des Charakters nicht überscheiten.

Wird der Zielende während des Zielens getroffen, so wird der angesammelte Zielen-Bonus entfernt.

\wbif{tirakan}{}{
\subsection{Angriffsmodi}
\index{Angriffsmodi}

Bei der jeweiligen Waffe ist angegeben, mit welchen \textit{Angriffsmodi} der Träger der Waffe sie verwenden kann. Der Spieler wählt für jeden Angriff beliebig aus den verfügbaren Modi. Das Umschalten des Feuermodus bei modernen Waffen erfordert keine Aktion.

\subsubsection{Nahkampf}
\index{Fertigkeit!Nahkampf}\index{Nahkampf}

Alle Nahkampfwaffen haben ausschließlich diesen Angriffsmodus. Der Charakter schlägt mit der Waffe im Nahkampf zu.
\begin{itemize}
    \item Der Angriff kann \textit{pariert} werden.
    \item Dem Angriff kann \textit{ausgewichen} werden.
    \item Der Angriff kann \textit{kritische Treffer} verursachen.
    \item Für den Angriff kann der Charakter zuvor \textit{zielen}.
\end{itemize}

\subsubsection{Einzelschuss}
\index{Einzelschuss}

Es wird ein Schuss pro Aktion abgegeben. Dies trifft auf viele moderne Waffen, aber auch auf Bögen, Schleudern und Armbrüste zu.
\begin{itemize}
    \item Der Angriff verbraucht 1 Munition.
    \item Der Angriff kann \textbf{nicht} \textit{pariert} werden.
    \item Dem Angriff kann \textbf{nicht} \textit{ausgewichen} werden.
    \item Der Angriff kann \textit{kritische Treffer} verursachen.
    \item Für den Angriff kann der Charakter zuvor \textit{zielen}.
\end{itemize}

\subsubsection{Feuerstoß}
\index{Feuerstoß}

Die Waffe wird im Feuerstoß Modus verwendet, es wird ein kurzer Feuerstoß abgegeben, welcher etwas ungenauer als ein Einzelschuss ist.
\begin{itemize}
    \item Dem Angriffswurf werden 2 Würfel hinzugefügt.
    \item Der Angriff verbraucht 3 Munition.
    \item Der Angriff kann \textbf{nicht} \textit{pariert} werden.
    \item Dem Angriff kann \textbf{nicht} \textit{ausgewichen} werden.
    \item Der Angriff kann \textbf{keine} \textit{kritischen Treffer} verursachen.
    \item Für den Angriff kann der Charakter \textbf{nicht} \textit{zielen}.
\end{itemize}
}

\subsection{Falsche Distanz}
\index{Falsche Distanz}

Jede Waffe hat eine angegebene Distanz, auf der sie effektiv ist. Weicht die Distanz des Ziels von der bei der Waffe angegebenen ab, ergibt sich ein Malus auf die Trefferwürfel.

Wenn die wirkliche Schussdistanz geringer ist als die angegebene Distanz der Waffe, wird der Angriff normal durchgeführt. Ist die Distanz bis zum Doppelten der Waffe erhöht, ist der Mindestwurf des Angriffs um 2 erhöht.

Wenn die Entfernung des Ziels mehr als das Doppelte der Waffenreichweite entfernt ist, ist es nicht möglich auf das Ziel zu schießen oder anzugreifen.

\subsection{Deckung}
\index{Deckung}

Sind Teile des Angegriffenen aus der Sicht des Angreifers verborgen, so gilt die Regel der Deckung. Hierbei kommt es darauf an, wie sehr der Angegriffene verborgen ist. Die Deckung ist hierbei in 3 Stufen eingeteilt:
\begin{itemize}
    \item 4+ Deckung: Das Meiste des Angegriffenen ist verborgen
    \item 5+ Deckung: Der Angegriffene ist halb verborgen
    \item 6+ Deckung: Es ist etwas schwerer den Angegriffenen hinter leichter Deckung zu treffen. Dieser Effekt wird u.A. durch den ``In der Hocke'' Zustand erreicht.
\end{itemize}

Hat der Angegriffene mindestens eine 6+ Deckung, so steht ihm nach dem \textit{Trefferwurf} ein Deckungswurf zu. Hierfür wirft er so viele Würfel, wie der Angreifer \textit{Treffer} hatte. Für jeden Erfolg (auf dem Mindestwurf gemäß der Deckung) wird ein Treffer entfernt.

Wenn der Angreifer \textit{kritische Treffer} erzielt hat, muss der Deckungswurf getrennt für kritische und normale Treffer gewürfelt werden, um festzustellen, welche Treffer verhindert wurden.

\subsubsection{Schilde}
\index{Schilde}

Schilde können verwendet werden, wenn der Charakter eine Einhandwaffe führt.

Schilde können auf zwei unterschiedliche Arten verwendet werden.
\begin{itemize}
    \item Für den \textbf{Schildblock} wird das Schild in der eigenen Runde mit zwei Aktionen bereit gemacht. In den folgenden Kampfrunden bietet das Schild die unten genannte Deckung für alle Angriffe gegen den Charakter. Während der Schildblock aktiv ist, ist die Bewegungsreichweite des Charakters halbiert. Der \textbf{Schildblock} ist so lange aktiv, bis der Charakter ihn aufhebt, also das Schild senkt.
    \item Die \textbf{Schildparade} kann spontan als \textit{Reaktion} verwendet werden. Sie bietet den u.g. Deckungswurf für einen einzelnen Angriff und kostet eine Aktion.
\end{itemize}

Schilde haben im Gegensatz zu anderen Rüstungen einen besonderen Wert, den Deckungswert. Dieser wird in der Form X+ angegeben, was bedeutet, dass Schilde Deckung in dieser Höhe bieten. Ein Rundschild bietet 5+ Deckung, nach einem Angriff kann der Angegriffene also für jeden Treffer 5+ würfeln, um ihn \textit{vor} der Anwendung von \textit{Schutz} und \textit{Wunden} zu verhindern. Dies ist sowohl mit \textit{Schildparade} als auch mit \textit{Schildblock} möglich.

\subsection{Schutz und Durchschlag}
\index{Schutz und Durchschlag}

Jeder Erfolg des \textit{Trefferwurfs}, welcher nicht durch \textit{Deckung} verhindert wurde, ist ein \textit{Treffer} beim Ziel des Angriffs. Auch andere Umstände können \textit{Treffer} verursachen, so kann eine Explosion z.~B. ``3 Treffer mit je 2 Wunden'' verursachen. Hier können Treffer durch Deckung verhindert werden.

Nimmt ein Charakter \textit{Treffer} hin, so kann er \textit{Schutz} einsetzen, um diese Treffer zu verhindern. Der Charakter hat einen \textit{Schutzpool}, welcher sich aus aller Rüstung und eventuellen anderen Effekten zusammensetzt. Für jede eingesetzte Schutzeinheit wird ein Treffer verhindert, eventuell mit weiteren Effekten (siehe Schutzpool).

Jeder nicht durch \textit{Schutz} verhinderte Treffer wird zu so vielen Wunden, wie es bei der Waffe oder dem Effekt angegeben ist. Ist nichts angegeben, so verursacht ein Treffer eine Wunde.

\subsection{Der Schutzpool}
\index{Schutzpool}

Jeder Charakter hat einen \textit{Schutzpool}, der sich aus allen Rüstungsteilen zusammensetzt. Jedes Rüstungsteil hat einen bestimmten Schutz, der in Form von Schutzeinheiten angegeben wird. Näheres zu den den Schutzeinheiten findet man im Kapitel \cref{ch:equipment}.

Wird ein Charakter angegriffen oder auf andere Weise getroffen, kann er Schutz aus seinem Schutzpool einsetzen, um diese Treffer zu verhindern. Der Einsatz von Schutz kostet keine Aktion und es kann beliebig viel Schutz eingesetzt werden.

Der Schutzpool stellt die Rüstung dar, die ein Charakter im Kampf trägt. Während des Kampfes verrutscht diese, vielleicht reißt ein Riemen und ein Teil der Rüstung fällt ab. So wird der Pool während des Kampfes kleiner, was durch das Ausgeben von Schutz dargestellt wird. Nach dem Kampf werden alle Schutzeinheiten im Pool wiederhergestellt.

Der Schutzpool steht nur im Kampf zur Verfügung. Erleidet ein Charakter außerhalb des Kampfes Treffer, so liegt es an ihm und der Spielleitung, die mögliche Schadensminderung durch Rüstung zu bewerten.

\wbif{tirakan}{
\begin{pexample}[Beispiel]
Hagen trägt einen Wappenrock aus Leinen. Dieser hat den Schutzwert ``N N B'', also zweimal normalen Schutz und einmal Schutz vor Blutung. Sein Gegner schlägt mit einer gewaltigen Kriegsaxt zu und verusacht 6 Treffer und einen kritischen Treffer.

Hagens Spieler entscheidet, dass er zwei der normalen Treffer mit dem ``N'' Schutz entfernt. Er streicht für diesen Kampf den N Schutz aus seinem Schutzpool und nimmt noch 4 Treffer hin. Die kritischen Treffer kann er nicht verhindern, da weder N (normaler Schutz) noch B (Schutz vor Blutung) kistische Treffer verhindern kann.
\end{pexample}
}{
\begin{pexample}[Beispiel]
John ist gerade aus dem Krankenhaus geflohen. Er trägt ein OP-Kleid mit einem ``B'' Schutz, also einem Schutz vor Blutungen. Auf der Strasse wird er mit einem Messer angegriffen. Der Angreifer wirft ``1 Crit, 1 Wunde''. John kann den kritischen Treffer nicht verhindern, verwendet jedoch den ``B'' Schutz um den normalen Treffer zu verhindern. Er streicht für diesen Kampf den ``B'' Schutz aus seinem Schutzpool und nimmt eine Wunde vom kritischen Treffer hin.
\end{pexample}
}

\subsection{Wunden}
\index{Wunden}

Eine \textit{Wunde} wird direkt zu den vom Charakter hingenommenen Wunden hinzugefügt. Sie kann nur verhindert werden, wenn eine Schablone, Ausrüstung oder anderes explizit eine Regel enthält, welche Wunden verändert.

\subsection{Waffenloser Nahkampf}
\index{Waffenloser Nahkampf}

Greift der Character ohne Waffe an, so wirft der Spieler Trefferwürfel entsprechend seinem Wert \textit{Nahkampf}. Der Mindestwurf entspricht dem Mindestwurf des Charakters, also in der Regel 5+.

Ist der Wert \textit{Kraft} des Charakters höher als 2, so hat der Nahkampfangriff \textit{Durchschlag} 1.

Ist der Wert \textit{Schnelligkeit} des Charakters höher als 2, so fügt der Charakter dem Wurf einen Würfel hinzu.

Die Reichweite eines waffenlosen Nahkampfangriffs beträgt 1 Meter.

\subsection{Ausweichen}
\index{Ausweichen}

Der Angegriffene kann einem Nahkampfangriff als Reaktion ausweichen. Voraussetzung hierfür ist, dass der angegriffene Charakter eine Aktion verfügbar hat und den Angreifer wahrnehmen kann. Einem Angriff von hinten kann also nicht ausgewichen werden.

Der Wert ist gleich dem Ausweichwert der Charakterschablonen plus dem Durchschnitt von Schnelligkeit und Geschicklichkeit (aufgerundet). Die Belastung durch Rüstung und Waffen verringert diesen Wert.

Um einem Angriff auszuweichen, wirft der Charakter einen Wurf auf seinen Wert in \textit{Ausweichen}. Der Mindestwurf hierfür ist um die Anzahl an Treffern des Gegners erhöht. Erzielt der Angegriffene mindestens einen Erfolg, so ist er dem Angriff komplett ausgewichen.

\section{Nahkampfangriffe parieren}
\index{Parade}\index{Nahkampf!Parade}

Nahkampfangriffe können pariert werden, wenn der Angegriffene eine passende Nahkampfwaffe bereithält und eine Aktion übrig hat.

Hierzu wird als \textit{Reaktion} wie auf einen Angriff mit der Waffe geworfen. Für jeden Erfolg bei diesem Wurf wird ein Treffer des Angreifers entfernt. \textit{Kritische Treffer} werden bei der Parade nur durch kritische Erfolge verhindert.

\section{Besondere Angriffe}
\index{Angriff!besondere Angriffe}

Es gibt eine Reihe besonderer Angriffe, mit denen ein Charakter seinen Angriff verfeinern oder verändern kann.

\subsection{Genauer Angriff}
\index{Genauer Angriff}

Beim genauen Angriff zielt der Charakter länger, um einen besseren Treffer zu landen. Das Tauschverhältnis ist dabei 1 Aktion für die Reduktion des Mindestwurfes um 1. Das Tauschen kann auch über Runden hinweg gehen. Der Mindestwurf kann maximal um den Wahrnehmungswert des Charakters verringert werden, kann aber 2 nicht unterschreiten. In der Zeit kann keine andere Aktion gemacht werden. Danach wird mit den veränderten Werten ein normaler Angriff gemacht.

\subsection{K.O. Angriff}
\index{K.O. Angriff}

Der K.O. Angriff hat nur die Absicht, einen Gegner K.O. zu schlagen, ohne ihm jedoch Schaden zuzufügen. Dabei muss der Angreifer eine stumpfe Waffe führen, oder zumindest mit einem stumpfen Gegenstand zuschlagen. Ist der Angriff erfolgreich, so würfelt der Gegner auf seine Resistenz. Erreicht er dabei nicht genauso viele Erfolge wie es Treffer gibt, geht er K.O.

Der Angriff macht keine Wunden. Deckung und Rüstung werden wie gewohnt berücksichtigt.

\subsection{Gewaltiger Angriff}
\index{Gewaltiger Angriff}

Bei einem gewaltigen Angriff nimmt der Charakter alle seine Kraft zusammen, um einen gewaltigen Schlag zu wirken. Pro zusätzlicher Aktion, die der Charakter in diesen Angriff investiert, erhöht sich die Anzahl der Würfel für diesen Angriff um 3, bis maximal zum Kraftwert des Charakters.

\subsection{Entwaffnender Angriff}
\index{Entwaffnender Angriff}

Mit einem entwaffnenden Angriff versucht der Angreifer, dem Gegner die Waffe aus der Hand zu schlagen. Dazu muss ihm ein Angriff auf den Waffenarm gelingen, dessen Mindestwurf um 2 angehoben ist. Der Angegriffene muss nach dem Angriff auf seine Kraft oder sein Geschick werfen, und mindestens so viele Erfolge erreichen wie der Angreifer Treffer hatte.

Gelingt dem Angegriffenen dies nicht, so wurde er entwaffnet.

Der entwaffnende Angriff verursacht keinen Schaden.

\subsection{Beidhändiges Kämpfen}
\index{Beidhändiges Kämpfen}

Ist der Charakter besonders geschickt im Umgang mit einer Waffe, so kann er zwei Waffen der gleichen Art gleichzeitig, also beidhändig führen. Das beidhändige Kämpfen ist nur mit Einhandwaffen möglich. Waffen, die sowieso mit beiden Händen geführt werden (Schwere Äxte, Stangenwaffen etc.) können nicht im beidhändigen Kampf geführt werden.

Führt ein Charakter zwei Waffen der gleichen Art gleichzeitig, erhält der Charakter eine Aktion mehr pro Kampfrunde. Die Waffe, die er mit der Sekundärhand führt, greift mit einem um 1 erhöhten Mindestwurf an.

\subsection{Waffe abstützen}
\index{Waffe abstützen}

Wenn dies mit der verwendeten Waffe möglich ist (in der Regel Schusswaffen ausser Bögen) kann der Charakter die Waffe vor der Benutzung an einer geeigneten Stelle auflegen. Das Abstützen dauert eine Aktion. Wird mit einer abgestützten Waffe geschossen, so ist der Mindestwurf um 1 verringert. Es kostet keine Aktion eine abgestützte Waffe wieder aufzunehmen.

\subsection{Coup de grâce}
\index{Coup de grâce}

Ein Charakter kann einen Gegner direkt töten, wenn dieser \textit{bewustlos}, \textit{schlafend} oder \textit{sterbend} ist. Hierzu wirft der Spieler einen normalen Angriffswurf. Ist dieser Wurf mit mindestens einem Erfolg gelungen, so erhält der Gegner den Status \textit{sterbend} mit der Stufe entsprechend den Erfolgen des Angriffs. Ist der Gegner bereits \textit{sterbend}, so wird die Stufe des Zustands um die Anzahl der Wunden des Angriffs erhöht.

Misslingt der Angriff, so wird ein schlafendes Opfer vermutlich erwachen.

\section{Werfen von Gegenständen}
\index{Werfen!Gegenstände}\index{Gegenstände!werfen}
\label{sec:throwing}

Wird ein Gegenstand, etwa eine Wurfnetz, auf ein Ziel geworfen, so wirft der Charakter auf seinen Wert \textit{Werfen}. Der Mindestwurf entspricht dem Mindestwurf des Charakters, in der Regel also 5+.

Ergibt der Wurf mindestens einen Erfolg, so hat der Charakter sein Ziel getroffen.

\subsection{Abweichung}
\index{Abweichung}

Zeigt der Wurf auf \textit{Werfen} keinen Erfolg, so ist der Wurf fehlgeschlagen. In diesem Fall wird auf die Abweichung geworfen.

Zunächst werden 2W6 geworfen, um die Richtung der Abweichung zu ermitteln. Das Resultat des Wurfes ergibt die Richtung in Form der ``Uhrzeit'', gesehen von dem werfenden Charakter mit Blick auf das Ziel. Eine 3 weicht somit, gesehen vom Werfenden, nach rechts vom Ziel ab.

Danach wird ein W6 geworfen, welcher die Entfernung der Abweichung in Metern ermittelt.

Das geworfene Objekt landet somit an dem ermittelten Ort.
\end{multicols}
